\scanpage{25}
\deleted{Soit en effet} Si on a $U_s\subset\deleted{\unsure{E|I}}\Sigma_i^-$, $\overline{U_s}\deleted{\unsure{=}}\cap E_{\unsure{i|j}}\ne\varnothing$, \deleted{et si} alors on a
\[
\begin{aligned}
U_s=X\setminus\mathop{\unsure{\rlap{\raisebox{.65ex}{$\scriptstyle-$}}\bigcup|\deleted{\bigcap}\bigcup}}\limits_{\alpha\in K}\widehat{\Sigma}_\alpha^+
&\qquad\text{où}\quad i\in K\subset I\\
&\forall\alpha\in K\qquad\Sigma_\alpha^+\in\pi_0(X\setminus E_\alpha)\ \deleted{\unsure{\text{les}|K}}\\
&\deleted{\unsure{U\Sigma|\bigcup\Sigma|U}}\ \text{les }\Sigma_\alpha^+\ \text{\unsure{mut\textsuperscript{t}.|mut.} disjoints}
\end{aligned}
\]
Il faut voir que $\deleted{K}\inserted{K\setminus\{i\}}$ et les $\Sigma_\alpha^+$ $(\alpha\in K\setminus\{i\})$ sont \unsure{determinés|déterminés|determines} en termes de $i$ et de $\Sigma_i^+$ i.e. $\Sigma_i^-$. Je dis que
\[
K\setminus\{i\}=\bigl\{\alpha\in I\mid E_\alpha\subset\Sigma_\alpha^-
\]
% archive sheet number: 24
